\section{结论}
典型日确定性模型在首末库存相同的条件下得到35126.95元购电费，较无储能节省26.90\%，展示了储能协调时段供需与电价的经济作用。历史场景日前模型的年度费用为14389135.08元，有效残差筛选带来的12861.89元改善主要集中于启动期。

预报融合与场景树滚动主方案的费用为14022396.98元，较其同口径历史预测对照减少367495.52元。6点与12点更新在本年度具有较明确的增量收益，18点在固定电价主模型中未表现出额外节省。波动电价下，4-2主模型实际费用15163043.89元，均价对照略低，说明价格与缺口关系的建模价值必须由具体回测评估。

4-3既有策略的实际费用为15061735.43元，与第三问的费用及调度差异已经分项呈现，但其优化目标存在额外购电价格项，故当前结果应作为具有该实现条件的策略输出保留。两种效率口径的已有对照进一步表明，物理损耗假设会改变费用数值。以上结论共同强调：比较购电策略前，应明确物理、信息与结算条件，并把可核算的结果与尚待验证的解释分开。
